\documentclass{article}

\usepackage{amsmath,amssymb}
\usepackage{xurl}
\usepackage[hidelinks]{hyperref}
\setlength{\emergencystretch}{2em}

\input{command}

\title{Wolf Proposition~6.2 --- Trivial Jordan Blocks for Peripheral Eigenvalues}
\author{}
\date{2026-08-05}

\begin{document}
\maketitle
\gapnote{scope-restriction}{resolved}

\section{The assertion}

Let $T:\MN{d}\to\MN{d}$ be a positive linear map that is either
trace-preserving or unital. Proposition~6.2 of Ref.~\cite{Wolf2012Quantum}
states that every eigenvalue $\lambda$ of $T$ with $|\lambda|=1$ has equal
geometric and algebraic multiplicities. Equivalently, every Jordan block
associated with $\lambda$ is one-dimensional.

The proof in Ref.~\cite{Wolf2012Quantum} uses two facts. First, positivity
together with trace preservation or unitality makes the iterates of $T$
uniformly bounded in operator norm. Second, powers of a nontrivial Jordan block
whose eigenvalue has modulus one grow without bound. These facts are
incompatible.

\section{Bounded orbits}

The bounded-orbit result is formalized under both alternatives in
Proposition~6.2 of Ref.~\cite{Wolf2012Quantum}.\footnote{Formal declarations:
    \leanid{IsPositiveMap.hasBoundedOrbits_of_tracePreserving} in
    \doclink{QICLean/Channel/FixedPoint/MeanErgodicProjection}, and
    \leanid{IsPositiveMap.hasBoundedOrbits_of_unital} in
    \doclink{QICLean/Channel/Peripheral/JordanBlocks}.}
Both proofs reduce an arbitrary matrix to the positive-semidefinite case using
the same lemma.\footnote{Formal declaration:
    \leanid{IsPositiveMap.hasBoundedOrbits_of_posSemidef_orbit_bounded} in
    \doclink{QICLean/Channel/FixedPoint/MeanErgodicProjection}.}

For a positive semidefinite input, trace preservation fixes the trace of every
iterate. Unitality instead bounds each iterate in the Loewner order by a scalar
multiple of the identity and then bounds its trace. The common scalar estimate
used in Ref.~\cite{Wolf2012Quantum} is
\begin{align}
    \tr\!\left(A\,T(B)\right)
        &\leq \lVert A\rVert_\infty\lVert B\rVert_\infty
            \tr\!\left(\Id\,T(\Id)\right)
    \label{eq:wolf_prop62_jordan_blocks_trace_bound_first}\\
        &=d\lVert A\rVert_\infty\lVert B\rVert_\infty.
    \label{eq:wolf_prop62_jordan_blocks_trace_bound_dimension}
\end{align}
Under trace preservation,
\begin{align}
    \tr\!\left(T(\Id)\right)
        &=\tr(\Id)
    \label{eq:wolf_prop62_jordan_blocks_trace_preserving_identity}\\
        &=d.
    \label{eq:wolf_prop62_jordan_blocks_trace_preserving_dimension}
\end{align}
Under unitality,
\begin{align}
    T(\Id)&=\Id,
    \label{eq:wolf_prop62_jordan_blocks_unital_identity}\\
    \tr\!\left(\Id\,T(\Id)\right)
        &=\tr(\Id)
    \label{eq:wolf_prop62_jordan_blocks_unital_trace}\\
        &=d.
    \label{eq:wolf_prop62_jordan_blocks_unital_dimension}
\end{align}
Thus (\ref{eq:wolf_prop62_jordan_blocks_trace_bound_dimension}) holds under
either hypothesis. The formal proof follows this direct argument and does not
pass through the adjoint map.

\section{The Jordan-block argument}

Suppose that $X,Y\in\MN{d}$ and $\lambda\in\C$ satisfy
\begin{align}
    (T-\lambda\Id)X&=Y,
    \label{eq:wolf_prop62_jordan_blocks_rank_two_first}\\
    (T-\lambda\Id)Y&=0,
    \label{eq:wolf_prop62_jordan_blocks_rank_two_second}\\
    Y&\neq0.
    \label{eq:wolf_prop62_jordan_blocks_rank_two_nonzero}
\end{align}
These relations describe a rank-two generalized eigenvector. Expanding
$T=\lambda\Id+(T-\lambda\Id)$ and using
(\ref{eq:wolf_prop62_jordan_blocks_rank_two_first}) and
(\ref{eq:wolf_prop62_jordan_blocks_rank_two_second}) gives
\begin{align}
    T^nX
        &=\lambda^nX+n\lambda^{n-1}Y.
    \label{eq:wolf_prop62_jordan_blocks_iterate_formula}
\end{align}
This identity is formalized by
\leanid{IsPositiveMap.pow_apply_rank_two_genEig}.

When $|\lambda|=1$, the reverse triangle inequality applied to
(\ref{eq:wolf_prop62_jordan_blocks_iterate_formula}) yields
\begin{align}
    \lVert T^nX\rVert
        &\geq n\lVert Y\rVert-\lVert X\rVert.
    \label{eq:wolf_prop62_jordan_blocks_linear_growth}
\end{align}
Because $Y\neq0$, the right-hand side of
(\ref{eq:wolf_prop62_jordan_blocks_linear_growth}) is unbounded as
$n\to\infty$. This contradicts the bounded-orbit theorem under either trace
preservation or unitality. Higher-rank generalized eigenvectors reduce to the
rank-two case by induction on the length of the generalized-eigenvector chain.

\section{Verdict}

Both alternatives in Proposition~6.2 of Ref.~\cite{Wolf2012Quantum} are
formalized. The shared binomial expansion is
\leanid{IsPositiveMap.pow_apply_rank_two_genEig}; the rank-two prohibitions and
the full generalized-eigenvector reductions are
\leanid{IsPositiveMap.no_rank_two_genEigenvector_of_tracePreserving},
\leanid{IsPositiveMap.no_rank_two_genEigenvector_of_unital},
\leanid{IsPositiveMap.peripheral_Jordan_trivial_of_tracePreserving}, and
\leanid{IsPositiveMap.peripheral_Jordan_trivial_of_unital}, all in
\doclink{QICLean/Channel/Peripheral/JordanBlocks}.

The unital case was formerly recorded as a scope restriction while a proof
through the trace adjoint was under consideration. That route is unnecessary:
the estimate in
(\ref{eq:wolf_prop62_jordan_blocks_trace_bound_first})--%
(\ref{eq:wolf_prop62_jordan_blocks_trace_bound_dimension}) applies directly
under either hypothesis. The unital and trace-preserving cases are therefore
proved in parallel. The work is tracked in
\href{https://github.com/LionSR/TNLean/issues/3984}{TNLean issue \#3984} and
closed in
\href{https://github.com/LionSR/TNLean/issues/5447}{TNLean issue \#5447}.

\bibliographystyle{alpha}
\bibliography{references}

\end{document}
